\chapter{Implementation}
\label{ch:implementation}

% CHAPTER ROLE (do not let this blur into Ch5 or Ch7):
%   Ch5 owns WHY the structure is shaped this way (argued architecture).
%   Ch7 owns HOW WELL it works (measurements against RQ1--RQ4).
%   Ch6 owns HOW the architecture becomes running code: realization + evidence.
% Organised as a session-walkthrough spine (see section.md). Each section lives
% in its own file:
%   01_structure    -> \cref{sec:impl-structure}
%   02_session      -> \cref{sec:impl-session}      (the spine)
%   03_highlights   -> \cref{sec:impl-highlights}
%   04_engineering  -> \cref{sec:impl-engineering}
%   06_summary      -> \cref{sec:impl-summary}

This chapter shows how the architecture of \cref{ch:design} becomes running code.
Where the previous chapter argued why the platform is shaped as it is, this one is concerned with how that shape is realised and with the evidence that it behaves as intended, stopping short of the quantitative evaluation that \cref{ch:evaluation} provides.
Rather than catalogue the source file by file, the chapter takes a single experiment session as its organising spine: it first locates the code, then follows one session through it from setup to a sealed record, and only then opens the individual mechanisms that the session exercises.
The emphasis throughout is on the realisations that carry the four concerns of \cref{sec:design-modules} and answer the research questions of \cref{sec:final-problem-statement}, illustrated with short excerpts of the actual repository code rather than exhaustive listings.

\Cref{sec:impl-structure} locates the code, mapping the two monorepo roots and the project structure that enforces the architecture's inward-pointing dependency rule, and records the concrete technology stack behind the system.
\Cref{sec:impl-session} is the spine: it follows one experiment session end to end, from the researcher's first command to the sealed record, across the two screens and two channels through which the platform is operated.
\Cref{sec:impl-highlights} then opens the mechanisms that walkthrough names only in passing, organised by the four separated concerns and the research questions they answer.
\Cref{sec:impl-engineering} steps back to the engineering practices that hold the two code roots together, from continuous integration to the documentation produced alongside the code.
\Cref{sec:impl-summary} draws the chapter together and hands off to the evaluation.

% \screenshotfig{path}: include a UI screenshot if the file is present, otherwise
% render a "pending" placeholder so the build stays green until the image is added.
% Screenshots live in Chapters/06_Implementation/screenshots/ (raster, from the user).
\newcommand{\screenshotfig}[2][\linewidth]{%
  \IfFileExists{#2}%
    {\includegraphics[width=#1]{#2}}%
    {\fbox{\parbox[c][3cm][c]{0.88\linewidth}{\centering\itshape Screenshot pending: place the referenced image file under \texttt{Chapters/06\_Implementation/screenshots/}.}}}%
}

% Code-listing styles (csharp, ts) and the IDE-style theme are defined globally in
% Setup/Settings.tex. Snippets are short, illustrative excerpts of real repository
% code (paths cited in each caption).

\section{Code Structure and Solution Layout}
\label{sec:impl-structure}
% Authoring brief for this section lives in section.md (§6.1).

\Cref{sec:design-style,sec:design-modules} described the architecture as concerns, ports, and a rule that dependencies point inward.
This chapter goes inside that description to show how it is realised in code, beginning here with where the code lives before \cref{sec:impl-session} follows a session through it.
The system is a monorepo with two independent roots, a .NET backend solution under \texttt{Backend/} and a Next.js front-end under \texttt{Frontend/}, with no shared workspace manifest binding them, so that each is built and versioned on its own.
Their folder layout is shown in \cref{fig:impl-tree} and their dependency structure in \cref{fig:impl-structure}.

\begin{figure}[htbp]
\footnotesize
\dirtree{%
.1 \texttt{reading-the-reader-monorepo/}.
.2 \texttt{Backend/src/} \DTcomment{.NET solution}.
.3 \texttt{core/}.
.4 \texttt{ReadingTheReader.core.Domain/} \DTcomment{domain entities}.
.4 \texttt{ReadingTheReader.core.Application/} \DTcomment{ports, services, snapshots}.
.5 \texttt{ApplicationContracts/Realtime/} \DTcomment{the four concern modules}.
.3 \texttt{infrastructure/}.
.4 \texttt{ReadingTheReader.TobiiEyetracker/} \DTcomment{Tobii adapter}.
.4 \texttt{ReadingTheReader.RealtimeMessenger/} \DTcomment{WebSocket broadcaster}.
.4 \texttt{ReadingTheReader.Realtime.Persistence/} \DTcomment{file store + checkpoints}.
.3 \texttt{ReadingTheReader.WebApi/} \DTcomment{transport, composition root}.
.2 \texttt{Frontend/src/} \DTcomment{Next.js application}.
.3 \texttt{app/} \DTcomment{routes only (sidebar groups)}.
.3 \texttt{modules/pages/} \DTcomment{feature UI}.
.3 \texttt{redux/}, \texttt{lib/} \DTcomment{state, API, gaze socket}.
.3 \texttt{components/}, \texttt{hooks/} \DTcomment{shared UI}.
}
\caption{Curated layout of the monorepo (selected folders only). The backend groups projects into \texttt{core} (the \texttt{Domain} entities and the \texttt{Application} with its ports), \texttt{infrastructure} (the adapters), and the \texttt{WebApi} composition root; the four concern modules of \cref{sec:design-modules} appear as folders beneath the application contracts. The front-end separates routing (\texttt{app/}) from feature UI, shared state and transport, and shared components.}
\label{fig:impl-tree}
\end{figure}

\begin{figure}[htbp]
\centering
\resizebox{\linewidth}{!}{%
\begin{tikzpicture}[
  web/.style={rectangle, rounded corners=3pt, draw=dtured, very thick, fill=dtured!10, align=center, text width=2.7cm, minimum height=0.95cm, inner sep=3pt, font=\scriptsize},
  infra/.style={rectangle, rounded corners=3pt, draw=black!55, thick, fill=grey!25, align=center, text width=2.2cm, minimum height=0.95cm, inner sep=3pt, font=\scriptsize},
  app/.style={rectangle, rounded corners=3pt, draw=navyblue, very thick, fill=blue!10, align=center, text width=2.6cm, minimum height=0.95cm, inner sep=3pt, font=\scriptsize},
  dom/.style={rectangle, rounded corners=3pt, draw=navyblue, very thick, fill=blue!18, align=center, text width=2.0cm, minimum height=0.8cm, inner sep=3pt, font=\scriptsize},
  fe/.style={rectangle, rounded corners=3pt, draw=navyblue, thick, fill=navyblue!8, align=center, text width=3.4cm, minimum height=1.0cm, inner sep=3pt, font=\scriptsize},
  dep/.style={->, thick, black!60},
  fedep/.style={->, semithick, black!45},
  ttl/.style={font=\small\bfseries, text=black!75},
  note/.style={font=\scriptsize\itshape, text=black!70, align=center},
]
% ---- panel titles ----
\node[ttl] at (-2.3, 6.3) {Backend solution (.NET)};
\node[ttl] at ( 5.0, 6.3) {Frontend (Next.js)};
% ---- divider ----
\draw[dashed, black!30] (2.9,-1.0) -- (2.9,6.0);
% ---- backend projects ----
\node[web]   (web) at (-4.7, 5.3) {{\bfseries WebApi}\\{\scriptsize transport + composition root}};
\node[infra] (per) at (-4.7, 3.5) {{\bfseries Realtime}\\{\scriptsize Persistence}};
\node[infra] (msg) at (-1.7, 3.5) {{\bfseries Realtime}\\{\scriptsize Messenger}};
\node[infra] (tob) at ( 1.4, 3.5) {{\bfseries Tobii}\\{\scriptsize Eyetracker}};
\node[app]   (app) at (-1.7, 1.7) {{\bfseries Application}\\{\scriptsize contracts, orchestration}};
\node[dom]   (dom) at (-1.7, 0.0) {{\bfseries Domain}};
% ---- backend references (A -> B: A depends on B) ----
\draw[dep] (app) -- (dom);
\draw[dep] (per) -- (app);
\draw[dep] (msg) -- (app);
\draw[dep] (tob) -- (app);
\draw[dep] (tob) -- (dom);
\draw[dep] (web) -- (per);
\draw[dep] (web) -- (msg);
\draw[dep] (web) -- (tob);
\draw[dep] (web) -- (app);
% ---- frontend layer stack ----
\node[fe] (feapp)   at (5.0, 5.1) {{\bfseries app/}\\{\scriptsize routes \& layouts (with-/without-sidebar)}};
\node[fe] (femod)   at (5.0, 3.6) {{\bfseries modules/pages}\\{\scriptsize feature UI}};
\node[fe] (festate) at (5.0, 2.1) {{\bfseries redux} \& {\bfseries lib}\\{\scriptsize state, RTK Query API, gaze socket}};
\node[fe] (feshare) at (5.0, 0.6) {{\bfseries components} \& {\bfseries hooks}\\{\scriptsize shared UI}};
\draw[fedep] (feapp)   -- (femod);
\draw[fedep] (femod)   -- (festate);
\draw[fedep] (festate) -- (feshare);
% ---- legend ----
\node[note, anchor=north west] at (-6.1,-1.2) {Arrow $A \rightarrow B$: project (or layer) $A$ references / depends on $B$. In the backend every reference points inward, towards the \textbf{Domain} core.};
\end{tikzpicture}%
}
\caption{Code structure of the two monorepo roots. In the backend solution every project reference points inward (\texttt{WebApi}, the composition root, depends on \texttt{Application} and all infrastructure; each infrastructure adapter and \texttt{Application} depend inward towards \texttt{Domain}), which is the dependency rule of \cref{sec:design-style} enforced by the build rather than merely intended. The front-end is a single Next.js application whose source separates thin routing (\texttt{app/}) from feature UI (\texttt{modules/pages}), shared state and transport (\texttt{redux}, \texttt{lib}), and shared components.}
\label{fig:impl-structure}
\end{figure}

The backend solution realises the concentric structure of \cref{sec:design-style} directly as projects.
A \texttt{Domain} project at the centre holds the basic entities, an \texttt{Application} project holds the use-case services, the snapshots, and the port interfaces, three infrastructure projects supply the adapters for the Tobii device, realtime messaging, and persistence, and a \texttt{WebApi} project provides the HTTP and WebSocket transport and composes the whole at startup.
Each conceptual layer of the architecture is thus a separate compilation unit, which is what allows the dependency rule to be enforced rather than merely intended.

That rule is visible in the project references, and because a project cannot compile against a reference it does not declare, the references are evidence rather than documentation.
The \texttt{Domain} project references nothing; the \texttt{Application} project references only \texttt{Domain}; and each infrastructure adapter references inward to \texttt{Application}, with the Tobii project additionally referencing \texttt{Domain}, so that every reference points towards the centre, as \cref{fig:impl-structure} shows.
The one project that references outward is \texttt{WebApi}, which depends on \texttt{Application} and on all three infrastructure projects at once.
This is deliberate, and it is the single exception the style permits: as the composition root, \texttt{WebApi} alone must know the concrete adapters in order to register them against the core's ports at startup, while the core itself never names them.

The front-end is a single Next.js application whose source mirrors the same discipline of separating routing from logic.
The \texttt{app/} directory contains only routes and layouts, organised into two route groups, a researcher-facing group with the operator sidebar and a participant-facing group without it, which are the two surfaces of the container view of \cref{sec:design-style}.
Feature UI lives in \texttt{modules/pages}, cross-feature state and the typed REST clients in \texttt{redux}, and shared helpers including the realtime gaze socket in \texttt{lib}, with common components and hooks factored into \texttt{components} and \texttt{hooks}.
Routing therefore stays thin and hands off to feature logic that is kept separately, in the same spirit as the backend's separation of transport from core.

Because there is no generated package shared between the two roots, the TypeScript types the front-end uses are hand-written mirrors of the backend records, and keeping the two aligned is a manual discipline rather than a compiler-checked one.
With the territory mapped, the next section puts it in motion: \cref{sec:impl-session} follows a single experiment session through these projects, from the researcher's first setup action to the sealed record left behind.

\section{Anatomy of an Experiment Session}
\label{sec:impl-session}
% Authoring brief for this section lives in section.md (§6.2). The spine of the chapter.
% UML class-diagram primitives (multipart class boxes + UML arrowheads) for this chapter's structural figures.
\usetikzlibrary{shapes.multipart, arrows.meta}

\Cref{sec:impl-structure} mapped where the code lives; this section sets it in motion by following one experiment session from the researcher's first preparation step to the sealed record it leaves behind.
It is the concrete counterpart of the abstract adaptive loop of \cref{sec:design-loop}, and it is the part of the system a researcher actually operates.

A single component, the \texttt{ExperimentSessionManager}, orchestrates the session and is the one place its state lives.
Its responsibilities are divided by concern across partial classes (a lifecycle facet, a gaze facet, and analysis, decision, intervention, and replay facets), but they share one object behind one lifecycle lock and reach the rest of the system only through ports.
\Cref{fig:impl-session-manager} shows that arrangement at the level of types: the manager is an \texttt{Application}-layer class that realises its segregated runtime facades and depends only on port interfaces also declared in \texttt{Application}, while the concrete adapters that drive the Tobii device, broadcast to the connected clients, and persist the session record live in the infrastructure projects and realise those ports.
Every realisation and dependency edge in the figure therefore crosses the \texttt{Application}/infrastructure boundary inward, which is the dependency rule of \cref{sec:design-style} holding at the granularity of individual classes rather than whole projects.
Every operation that changes session state then follows the same shape: it acquires the lifecycle lock, mutates the state, writes a checkpoint, and broadcasts the authoritative session snapshot that every connected surface renders.
Because all session state lives behind this one component and is published as a single snapshot, the researcher's control surface and the participant's reader never hold competing copies of the truth.
The session moves through three phases, configuring, running, and completed, with a reset back to the beginning, as \cref{fig:impl-session-states} shows.

\begin{figure}[htbp]
\centering
\resizebox{\linewidth}{!}{%
\begin{tikzpicture}[
  font=\scriptsize,
  mgr/.style={rectangle split, rectangle split parts=3, draw=navyblue, very thick, fill=blue!8, align=left, text width=5.0cm, inner sep=4pt},
  iface/.style={rectangle split, rectangle split parts=2, draw=black!65, thick, fill=grey!15, align=left, text width=3.7cm, inner sep=3pt},
  adapter/.style={rectangle split, rectangle split parts=2, draw=dtured, very thick, fill=dtured!8, align=left, text width=3.7cm, inner sep=3pt},
  realise/.style={dashed, -{Triangle[open, length=7pt, width=7pt]}, semithick, black!70},
  use/.style={dashed, -{Stealth[length=5pt]}, semithick, black!55},
  bnd/.style={dashed, black!40, thick},
  note/.style={font=\scriptsize\itshape, text=black!75},
]
% ---- core: the session manager (Application layer) ----
\node[mgr] (mgr) at (0,5.9)
  {{\bfseries ExperimentSessionManager}\\{\itshape sealed, partial}\\[2pt]{\tiny realises \mbox{IExperimentRuntimeAuthority}, \mbox{IExperimentSessionManager},\\\mbox{IExperimentSessionQueryService}, \mbox{IExperimentReplayRecoveryBuffer}}
   \nodepart{two}$-$\_lifecycleGate : SemaphoreSlim\\$-$\_session : ExperimentSession
   \nodepart{three}$+$StartSessionAsync() : bool\\$+$FinishSessionAsync(cmd) : Snapshot\\$+$GetCurrentSnapshot() : Snapshot};
% ---- core: the port interfaces (Application layer) ----
\node[iface] (psen) at (-5.8,1.5)
  {{\itshape interface}\\{\bfseries IEyeTrackerAdapter}
   \nodepart{two}$+$GazeDataReceived : event\\$+$StartEyeTracking()\\$+$BeginCalibrationAsync()};
\node[iface] (pbrd) at (0,1.5)
  {{\itshape interface}\\{\bfseries IClientBroadcaster\\Adapter}
   \nodepart{two}$+$BroadcastAsync(type, \mbox{payload})\\$+$ConnectedClients : int};
\node[iface] (psto) at (5.8,1.5)
  {{\itshape interface}\\{\bfseries IExperimentStateStore\\Adapter}
   \nodepart{two}$+$SaveActiveReplayAsync(doc)\\$+$LoadActiveReplayAsync() : Replay?};
% ---- Application / infrastructure boundary ----
\draw[bnd] (-7.8,-0.6) -- (7.8,-0.6);
\node[note, anchor=west] at (8.1,-0.12) {Application};
\node[note, anchor=west] at (8.1,-0.52) {(core, ports)};
\node[note, anchor=west] at (8.1,-1.0) {Infrastructure};
\node[note, anchor=west] at (8.1,-1.4) {(adapters)};
% ---- infrastructure: concrete adapters ----
\node[adapter] (asen) at (-5.8,-2.9)
  {{\bfseries TobiiEyeTracker\\Adapter}\nodepart{two}{\itshape project:\\TobiiEyetracker}};
\node[adapter] (abrd) at (0,-2.9)
  {{\bfseries WebSocketRealtime\\Messenger}\nodepart{two}{\itshape project:\\RealtimeMessenger}};
\node[adapter] (asto) at (5.8,-2.9)
  {{\bfseries FileSnapshotExperiment\\StateStoreAdapter}\nodepart{two}{\itshape project:\\Realtime.Persistence}};
% ---- dependencies: the manager uses the ports ----
\draw[use] (mgr.south) -- (pbrd.north) node[note, midway, fill=white, inner sep=1pt] {depends on};
\draw[use] (mgr.south west) to[out=-120,in=65] (psen.north);
\draw[use] (mgr.south east) to[out=-60,in=115] (psto.north);
% ---- realisations: the adapters realise the ports ----
\draw[realise] (asen.north) -- (psen.south) node[note, midway, fill=white, inner sep=1pt] {realises};
\draw[realise] (abrd.north) -- (pbrd.south);
\draw[realise] (asto.north) -- (psto.south);
\end{tikzpicture}%
}
\caption[Ports-and-adapters realisation around the session manager.]{Class view of the experiment session as ports and adapters. The \texttt{ExperimentSessionManager} (top) is an \texttt{Application}-layer class that realises its segregated runtime facades and depends only on the three port interfaces in the middle row (sensing, broadcast, and persistence). Each concrete adapter in the bottom row lives in an infrastructure project and realises the port above it. The dashed rule marks the \texttt{Application}/infrastructure boundary: every realisation (hollow triangle) and dependency (filled arrow) edge crosses it inward towards the core, so the dependency rule of \cref{sec:design-style} is visible at the level of individual classes. Members are abbreviated; the replay, recovery, and module-provider ports follow the same pattern and are omitted for clarity.}
\label{fig:impl-session-manager}
\end{figure}

\begin{figure}[htbp]
\centering
\resizebox{0.97\linewidth}{!}{%
\begin{tikzpicture}[
  state/.style={rectangle, rounded corners=4pt, draw=navyblue, very thick, fill=blue!8, align=center, text width=3.7cm, minimum height=1.5cm, inner sep=4pt, font=\small},
  lbl/.style={font=\scriptsize\itshape, text=black!70, align=center, fill=white, inner sep=1.5pt},
  arr/.style={->, thick, black!65},
]
\fill[black!75] (-9.0,0) circle (2.4pt);
\node[state] (cfg)  at (-6.6,0) {{\bfseries Configuring}\\[2pt]{\scriptsize gates: participant, eye tracker + licence, calibration + validation, reading material}};
\node[state] (run)  at ( 0.0,0) {{\bfseries Running}\\[2pt]{\scriptsize per item: read with the adaptive loop, then comprehension quiz}};
\node[state] (done) at ( 6.6,0) {{\bfseries Completed}\\[2pt]{\scriptsize record sealed and exported}};
\draw[arr] (-9.0,0) -- (cfg.west);
\draw[arr] (cfg.east) -- (run.west) node[lbl, midway, above] {start\\\mbox{[all gates ready]}};
\draw[arr] (run.east) -- (done.west) node[lbl, midway, above] {finish};
\draw[arr] (run.north) to[out=60,in=120,looseness=8] (run.north);
\node[lbl] at (0,1.75) {next item};
\draw[arr] (done.south) to[out=-118,in=-62,looseness=0.9] (cfg.south);
\node[lbl] at (0,-2.15) {reset};
\end{tikzpicture}%
}
\caption{Lifecycle of an experiment session. The session is configured behind four readiness gates, runs item by item once all gates are satisfied, and is sealed and exported on completion; a reset returns an inactive session to the configuring phase. Every transition passes through the single \texttt{ExperimentSessionManager}, which checkpoints the state and broadcasts the authoritative snapshot.}
\label{fig:impl-session-states}
\end{figure}

In the configuring phase the researcher prepares the session through a guided, role-sequenced stepper, shown in \cref{fig:impl-setup}.
The backend tracks four independent readiness gates, a participant, a selected eye tracker with a valid licence, a passed calibration, and the reading material, each of which reports whether it is satisfied and, if not, the reason it blocks the start.
The stepper sequences these into a recommended order, researcher preparation (the eye tracker and the reading material) before participant setup (the participant details and calibration), and locks each step until its prerequisites are met, so the researcher is never offered a control that cannot yet be used.
The eye-tracker step in the figure shows the detail of this preparation: a real Tobii device is selected and its licence is managed, with an explicit overwrite guard so that a saved licence is not replaced by accident.

\begin{figure}[htbp]
\centering
\screenshotfig[0.95\linewidth]{Chapters/06_Implementation/screenshots/Stepper.png}
\caption{The setup stepper in the configuring phase. The session is prepared through role-sequenced steps (researcher preparation before participant setup); each step is locked until its prerequisites are met, and the current step states the reason it still blocks the session start. Shown here is the eye-tracker step, with a real Tobii device selected and its licence guarded against accidental overwrite.}
\label{fig:impl-setup}
\end{figure}

Calibration is one of the four gates and runs as its own immersive sub-flow, shown in \cref{fig:impl-calibration}.
The participant follows a sequence of on-screen targets (sixteen in the configuration shown) while the device collects samples, after which a calibration is computed and applied and a validation pass measures its accuracy and precision.
Only a calibration whose validation passes satisfies the gate, and the accuracy and precision it reports are carried in the setup snapshot and evaluated in \cref{ch:evaluation}.

\begin{figure}[htbp]
\centering
\screenshotfig[0.78\linewidth]{Chapters/06_Implementation/screenshots/Calibration.png}
\caption{The calibration sub-flow. The participant fixates a sequence of on-screen targets (sixteen here) while the device collects samples; the resulting calibration is then validated for accuracy and precision, and only a passing validation satisfies the calibration gate.}
\label{fig:impl-calibration}
\end{figure}

When all four gates are satisfied the researcher starts the session, and the manager re-checks readiness before committing to the run (\texttt{StartSessionAsync}).
If anything is incomplete it refuses with the blocking step's reason, so a misconfigured run cannot be started rather than merely being discouraged; otherwise it resolves the order of the experiment's material items (fixed or randomised), resets the gaze counters and the event sequence, and marks the session active.
From this point the adaptive loop of \cref{sec:design-loop} runs over the modules resolved from the registries of \cref{sec:design-modules}, and the researcher observes and steers it through the live control surface of \cref{fig:impl-live}.
That surface is the second screen of \cref{sec:design-style}: it mirrors the participant's text, reports the live end-to-end latency, and offers the manual interventions and typography controls through which the researcher adapts the presentation.
Its behaviour makes the commit discipline of driver~D7 visible in the interface, since appearance changes are applied immediately whereas typography changes wait for the selected boundary, which is the two-phase, context-preserving commit of \cref{sec:design-loop} exposed as an operator control.

\begin{figure}[htbp]
\centering
\screenshotfig[0.95\linewidth]{Chapters/06_Implementation/screenshots/ResearcherLiveView.png}
\caption{The researcher's live control surface during a running session. It mirrors the participant's reader (the second screen of \cref{sec:design-style}), reports the live end-to-end latency, and exposes the manual interventions and typography controls. Appearance changes apply immediately while typography changes wait for the selected commit boundary, which is the context-preserving discipline of driver~D7 surfaced as an operator control.}
\label{fig:impl-live}
\end{figure}

The session then proceeds item by item.
For each material the participant reads while the loop adapts the text, and afterwards answers a short comprehension quiz, and both the reading events and the answers are written into the session record alongside every gaze sample, intervention, and lifecycle event.

Finishing the session ends it and seals its record.
The manager stops the loop, flushes any buffered events, and builds two exports from the recovery buffer (\texttt{FinishSessionAsync}), a full raw export and a processed one, saves them as the latest, marks the recovery completed, and clears the active record, after which the session is sealed exactly as \cref{sec:design-data} described.
A sealed session is then available for replay, shown in \cref{fig:impl-replay}: the recording is reconstructed as a timed playback, with the gaze token heat map drawn over the text and a filterable timeline of lifecycle, quiz, and intervention events beside it.
Replay is therefore the operational proof of \cref{sec:design-data}, that a session can be re-presented and re-analysed from its record alone.

\begin{figure}[htbp]
\centering
\screenshotfig[0.95\linewidth]{Chapters/06_Implementation/screenshots/Replay.png}
\caption{Replay of a sealed session. The recording is reconstructed as a timed playback (left), with the gaze token heat map drawn over the reconstructed text and a filterable timeline of lifecycle, quiz, and intervention events (right). That a session replays from its record alone is the operational evidence of the reproducibility argued in \cref{sec:design-data}.}
\label{fig:impl-replay}
\end{figure}

The walkthrough has named several mechanisms only in passing: the lock-free gaze path that keeps the live latency low, the registries that resolve the modules the loop drives, and the export that replay consumes.
\Cref{sec:impl-highlights} opens these three for a closer look.

\section{Implementation Highlights}
\label{sec:impl-highlights}
% Authoring brief for this section lives in section.md (§6.3).
\todo[inline]{\S6.3 to be written. 2--4 deeper dives the walkthrough flags: real-time gaze hot path (lock-free, WebSocket; RQ2); pluggable modules + a worked "add an intervention" (RQ1); external-provider framework realized; session record schema + replay (RQ4). Each: prose + at most one short listing. Pick the dives during scaffolding.}

\section{Engineering Practice}
\label{sec:impl-engineering}
% Authoring brief for this section lives in section.md (§6.5).
\todo[inline]{\S6.5 to be written. Build \& CI (frontend-ci, backend-ci; per-subproject deps; Windows/Tobii constraint); cross-cutting concerns (error handling conventions, hot-path vs gated concurrency, configuration, logging style); testing approach; challenges \& resolutions (gaze-to-content mapping over Markdown, latency, Windows-only Tobii SDK, context preservation across re-render).}

\subsection{Documentation as an engineering deliverable}
\label{subsec:impl-docs-site}

We treat documentation as part of the engineering work of this thesis rather than as an afterthought to it.
The reasoning is intrinsic to the system's own claims: a platform whose central promise is that others can extend it (RQ1) is only as extensible as its contracts are documented, and a researcher-operated platform is only reproducible if its setup and operation are written down for someone other than its authors to follow.
We therefore produced a companion documentation site as a deliverable in its own right, developed alongside the code and in the same repository, published at \url{https://mrsachin7.github.io/reading-the-reader-monorepo/}.
The site is written for the developer and operator rather than the participant: alongside the architecture and data-flow overviews it provides a local-setup and run-an-experiment guide, a reference for the REST endpoints and the session-record export format, and, most relevant to the extensibility claim, the module-provider integration protocol together with a worked guide to adding a new analysis or decision module and mock providers to develop against.
Keeping it aligned with the implementation and mirroring the structure of this report is a deliberate practice: it turns the extension path argued in \cref{sec:design-extensibility} into something operational rather than implicit, and leaves the teams who continue the project a maintained entry point rather than a codebase to reverse-engineer.

\section{Summary}
\label{sec:impl-summary}

This chapter has shown how the architecture of \cref{ch:design} is realised in running code.
\Cref{sec:impl-structure} located that code in a monorepo of two independently built roots, a .NET backend whose project references enforce the inward-pointing dependency rule and a Next.js front-end that separates thin routing from feature logic, and recorded the concrete stack behind them.
\Cref{sec:impl-session} then set that structure in motion, following one experiment session from the researcher's first command to the sealed record: two browser surfaces over one front-end, two channels (REST commands and a realtime WebSocket stream), and a single \texttt{ExperimentSessionManager} that owns the session state behind a lifecycle gate and publishes one authoritative snapshot to both screens.

\Cref{sec:impl-highlights} then opened the mechanisms that walkthrough had named only in passing, organised by the four concerns the architecture separates and the research questions they answer: the registry-backed module seams and the external-provider protocol that make the loop pluggable (RQ1); the dwell-based, area-of-interest sensing pipeline that turns a raw gaze stream into oculomotor events on a lock-free path and survives reflow (RQ2); the capture-and-restore mechanism that preserves the reader's position across an intervention and grades its cost (RQ3); and the experiment templates, the versioned dual-format record, and the deterministic replay that make a session repeatable and auditable after the fact (RQ4).
\Cref{sec:impl-engineering} closed the chapter with the practices that hold the two roots together: their continuous-integration pipelines, the conventions for error handling, logging, configuration, and concurrency, the testing approach and its honest asymmetry between the tiers, and the documentation produced as a deliverable in its own right.

What the chapter has established is that the system exists and behaves as designed; what it has not done is measure how well it does so.
\Cref{ch:evaluation} takes that step, verifying the realised platform against its functional requirements and non-functional budgets and reporting the latency and calibration figures that this chapter has so far only promised.

